\section{CONCLUSIONS}
\label{conclusion}

%In this paper, we proposed a modeling and control framework to handle aerodynamic forces with a complex-shaped VTOL system, the jet-powered humanoid robot iRonCub.  First, a simplified robot model is designed to generate a dataset of aerodynamic forces for a given flight envelope, using CFD simulations. Then, an identification procedure is applied to identify the Aerodynamics Characteristics of the system, and the coefficients are used to design a model of the aerodynamic forces acting on the robot. 


In this paper, analytical models of aerodynamic forces acting on the jet-powered  humanoid robot iRonCub are identified using data coming from CFD analysis. The identified models are then used to design appropriate controllers for the robot flying in the given flight envelope. %Gain-scheduling and feedback linearization techniques are applied to improve the current controller towards aerodynamic effects in critical outdoor environment which have been tested in simulations showing better performances in tracking of a desired CoM position and guaranteeing stabilization of a reference trajectory.
Two techniques, \emph{gain-scheduling} and \emph{feedback linearization}, are applied to improve the flight controller in presence of aerodynamic effects. The controller was tested successfully through simulations in a critical outdoor environment.

Nevertheless, the proposed framework still requires several assumptions, both for performing the CFD analysis and in the controller design. Future work may consider to distribute the aerodynamic forces all over the major robot links for more realistic simulations, and improve the geometry of the iRonCub model for CFD analysis. Finally, further effort in control design is desired to improve the robustness of the controller in different flight envelopes.

%Future work may consider to distribute the aerodynamic forces all over the major links to be close to the reality and overcome the limits of using a more complex iRonCub model in CFD simulations. Finally, further effort in control design is desired to improve the robustness of the controller for a practical implementation and experiments on the real robotic platform iRonCub with aerodynamic effects considered.

